\section{Experimental Setup}
Evaluation choices for this problem are dominated by one fact: entanglements are rare, contiguous events, and the 25 recordings differ in gait, terrain contact, and which leg is affected. A window-level random split would place windows from the same recording on both sides of the partition, and because neighboring windows share almost all of their samples, that arrangement leaks information and inflates every score. All protocols below therefore group by whole recording so that no recording contributes to both training and test.

\subsection{Data splits}
We report two complementary protocols. The first is a fixed leakage-safe split, grouped at the recording level, of 16 training / 4 validation / 5 test recordings; the five test recordings (four positive, one negative) are held out entirely and used only for the final numbers. Normalization statistics and the Walking reference distribution used by the physics severity index are fit on training rows alone, so no test information reaches the model or the feature pipeline.

The fixed split gives a single, interpretable operating-point result, but with 15 positive recordings any one test partition samples only a handful of events, and the affected legs are unevenly represented (FR appears in only three positive recordings). To measure generalization across recordings rather than across an arbitrary partition, we adopt leave-one-recording-out (LORO) cross-validation over the 15 positive recordings as the headline protocol: each fold holds out one positive recording, trains on the rest (with the negatives available for training), and evaluates on the held-out recording, and we report the mean and standard deviation of the per-fold $F_1$. LORO is the fairer test of whether the detector transfers to an event it has never seen, and its variance across folds is itself informative.

\subsection{Dense evaluation and metrics}
Training uses a stride of 25 samples between windows for efficiency, but evaluation is dense: we slide the 200-sample window with hop 1 so that a decision is produced at every timestep, matching how the runtime consumes \texttt{/lowstate}. This is the strictest setting, since it scores the boundaries of an event where the window straddles the transition into and out of the entangled phase.

We use standard definitions, each computed on the dense windows. Precision is the fraction of flagged windows that are truly entangled, recall the fraction of entangled windows that are flagged, and $F_1$ their harmonic mean. PR-AUC is the area under the precision--recall curve and ROC-AUC the area under the receiver-operating curve, both threshold-free summaries of the binary detection head; PR-AUC is the more meaningful of the two here because the positive class is rare. For the per-leg attribution task we report multi-label exact-match: a window counts as correct only when the predicted set of entangled legs equals the true set exactly, which is a deliberately unforgiving criterion.

\subsection{Rule-based baseline}
Snagged legs resist downward thigh motion, and a natural non-learned detector fires when the resistive thigh torque of any leg is high. We implement this thigh-torque heuristic and evaluate it under the identical dense, grouped protocol, on the same windows and the same test recordings, so the comparison in Table~\ref{tab:detection} isolates the value of the learned model rather than any difference in how the data are prepared.

\subsection{Feature ablation}
The shipped model reads 60 channels: 48 raw signals and 12 engineered per-leg features. To attribute performance to feature groups rather than infer it, we retrain the full model from scratch on channel subsets and re-run LORO for each: engineered features only (12), raw signals only (48), raw plus engineered without the foot-contact channels (56), raw plus engineered without the IMU channels (52), and the full 60-channel set. Because each variant is retrained under the same protocol, Table~\ref{tab:ablation} reflects the marginal contribution of each group under LORO variance rather than a post-hoc feature-importance estimate.

\subsection{Calibration}
The detector emits a probability, and downstream thresholding and debouncing behave well only if that probability is meaningful. We assess calibration in three ways. Brier score measures the mean squared error of the probabilities against the binary outcomes, expected calibration error (ECE) measures the gap between confidence and accuracy across probability bins, and we additionally inspect saturation, the fraction of raw probabilities pinned near 1.0. Saturation is the practical motivation here: when a large share of probabilities sit at the ceiling, the false-alarm-rate threshold is forced to an extreme value, and temperature scaling \cite{guo2017calibration,niculescu2005predicting} is applied to de-saturate the output so that a usable operating threshold exists. We report Brier, ECE, and saturation before and after scaling, without assuming that de-saturation implies better reliability.

\subsection{Latency}
Because the detector must keep pace with the 500 Hz feature stream, we benchmark inference time. Latency is measured on the deployment ONNX Runtime \cite{onnxruntime}, pinned to a single CPU thread, timing the per-window forward pass together with the post-processing (sigmoid, temperature scaling, threshold, and debounce logic) rather than the network alone. We report mean, median, and 95th-percentile per-window latency against the 2.0 ms budget implied by the sample step. This is a benchmark of the deployment runtime on a development CPU; it is not measured on the Go2's onboard computer, and we present it as such.